\chapter*{APPENDIX A: SOURCE CODE}
\addcontentsline{toc}{chapter}{APPENDIX A: SOURCE CODE}

The complete source code is available on GitHub:

\begin{center}
\textbf{https://github.com/The-Token-Trio/125970-126690-126687-Assignment2-PLC}
\end{center}

The project can be run from the command line or the desktop UI:

\begin{lstlisting}
python3 main.py
python3 ui.py
python3 -m unittest discover -s tests -v
\end{lstlisting}

No external Python dependencies are required.

% Compact style for appendix listings: tiny font, no line numbers, tight spacing
\lstdefinestyle{appendix}{
  basicstyle=\footnotesize\ttfamily,
  numbers=left,
  numberstyle=\tiny\color{gray},
  numbersep=8pt,
  breaklines=true,
  lineskip=0.4ex,
  frame=single,
  backgroundcolor=\color{codebg},
  rulecolor=\color{codeframe},
  showstringspaces=false,
  tabsize=4,
  captionpos=b,
  inputencoding=utf8,
  extendedchars=true,
  literate={—}{{---}}1 {→}{{$\rightarrow$}}1
}

% -----------------------------------------------------------------------
\section*{A.1\quad Entry Points}
% -----------------------------------------------------------------------

\lstinputlisting[style=appendix,
  caption={main.py --- command-line entry point},
  label={lst:main}
]{../main.py}

\lstinputlisting[style=appendix,
  caption={ui.py --- Tkinter desktop workbench},
  label={lst:ui}
]{../ui.py}

% -----------------------------------------------------------------------
\section*{A.2\quad Core Components}
% -----------------------------------------------------------------------

\lstinputlisting[style=appendix,
  caption={components/tokens.py --- token type enumeration and Token dataclass},
  label={lst:tokens}
]{../components/tokens.py}

\lstinputlisting[style=appendix,
  caption={components/lexica.py --- lexical analyser},
  label={lst:lexica}
]{../components/lexica.py}

\lstinputlisting[style=appendix,
  caption={components/ast\_nodes.py --- AST node definitions},
  label={lst:astnodes}
]{../components/ast_nodes.py}

\lstinputlisting[style=appendix,
  caption={components/parser.py --- recursive-descent parser},
  label={lst:parser}
]{../components/parser.py}

\lstinputlisting[style=appendix,
  caption={components/symbol\_table.py --- scoped symbol table},
  label={lst:symboltable}
]{../components/symbol_table.py}

\lstinputlisting[style=appendix,
  caption={components/type\_checker.py --- static type checker with inference},
  label={lst:typechecker}
]{../components/type_checker.py}

\lstinputlisting[style=appendix,
  caption={components/interpreter.py --- tree-walking interpreter},
  label={lst:interpreter}
]{../components/interpreter.py}

\lstinputlisting[style=appendix,
  caption={components/pipeline.py --- shared pipeline (CLI, UI, and tests)},
  label={lst:pipeline}
]{../components/pipeline.py}

% -----------------------------------------------------------------------
\section*{A.3\quad Test Suite}
% -----------------------------------------------------------------------

\lstinputlisting[style=appendix,
  caption={tests/test\_pipeline.py --- automated regression suite (54 tests)},
  label={lst:tests}
]{../tests/test_pipeline.py}
